
为了避免在导入节文件时频繁输入复杂的文件路径, 用户需预先设定\newconcept{节文件路径前缀}{Section Path Prefix}, 当导入节文件或添加节链接时, 只需提供节文件路径中除路径前缀和文件后缀名以外的部分（即\newconcept{节路径}{Section Path}）即可. 节路径中去掉节文件名称的部分称作该节所在的\newconcept{群组路径}{Group Path}. 请注意, \textcolor{red}{\SimpleSystemTeX 提供的任何命令均无需输入完整的节文件路径}, 完整的节文件路径会导致群组结构无法正确解析.

系统默认以 \texttt{./sec/} 作为节文件路径前缀, 因此请确保你的项目符合以下结构:

\medskip\dirtree{%
    .1 {\meta{Project}/} .
    .2 {\meta{Project}.tex} \DTcomment{主文件} .
    .2 {simplesystemtex.sty} \DTcomment{SimpleSystemTeX宏包} .
    .2 {sec/} .
    .3 {\meta{Group}/} {\DTcomment{一级群组}} .
    .4 {\meta{Subgroup}/} {\DTcomment{二级群组}} .
    .5 {\meta{Section1}.tex} .
    .4 {\meta{Section2}.tex} .
}\medskip

以上面所示的文件结构为例, \meta{Section1} 的各个路径分别为:

\List{Section1的路径}(){
    - 节文件路径: \texttt{./sec/\meta{Group}/\meta{Subgroup}/\meta{Section1}.tex}
    - 节路径: \texttt{\meta{Group}/\meta{Subgroup}/\meta{Section1}}
    - 群组路径: \texttt{\meta{Group}/\meta{Subgroup}/}
}

你可以通过下面的命令来修改节文件路径前缀, 以适应已有的项目结构:

\UserCommand{修改节文件路径前缀}<\cs{SetSectionPath}>{
    \cs{SetSectionPath}
    \meta{界定符}\meta{路径前缀}\meta{界定符}
}[
    - \meta{路径前缀} : 新的节文件路径前缀, 必须以 \texttt{./}（当前目录）或 \texttt{../}（上级目录）开头, 并且以 \texttt{/} 结尾, 路径以主文件所在目录为基准点, 各层级间通过 \texttt{/} 分隔.
]

你可以使用成对的 \texttt{\{} 和 \texttt{\}} , 或其它任意两个相同的字符作为界定符, 请务必确保路径前缀中不包含界定符.

修改节文件路径前缀后, \SimpleSystemTeX 会以新的方式解析群组结构, 例如将路径前缀修改为 \texttt{./sec/\meta{Group}/} 后, \meta{Section1} 所在的一级群组将变为 \meta{Subgroup} , 请务必确保新的节文件路径前缀与预期的文件结构相匹配.

在主文件的正文区使用下面的命令来导入节文件:

\UserCommand{导入节文件}<\cs{ImportSection}>{
    \cs{ImportSection}
    \{\meta{节路径}\}
    (\meta{节显示名称})
    <\meta{自定义参数}>
}[
    - \meta{节路径} : 要导入的节路径, 其中不能包含特殊字符 \texttt{\cs{}} , \texttt{\%} , \texttt{\{} , \texttt{\}} 和连续的空格.
][
    - \meta{节显示名称} : 用于显示的节名称, 若此参数为空则默认为节文件名称.
    - \meta{自定义参数} : 在\seclink{导言区/样式设定}中使用该参数来实现各种复杂的排版逻辑.
]

当节文件名称中包含不可打印的特殊字符时, 你\textcolor{red}{必须提供经过转义的节显示名称}以确保节名称能够在目录、节标题和交叉引用中正确显示.

\Example{导入节文件-示例}{
    \par\noindent \cs{ImportSection} \{线性空间/Hamel基\}
    \par\noindent \cs{ImportSection} \{拓扑空间/内部\&闭包\}(内部\cs{\&} 闭包)
    \par\noindent \cs{ImportSection} \{测度空间/Lp收敛\}(\$L\^{}p\$ 收敛)
}

\SimpleSystemTeX 将根据节文件的导入顺序确定正文中各节的呈现顺序, 以及它们在目录中的条目顺序. 你也可以添加分块来对节进行分组. 请注意, 分块完全依赖于节文件导入顺序, 而与各节所处的群组结构无关.

\UserCommand{创建分块}<\cs{AddPart}>{
    \cs{AddPart}
    \{\meta{分块}\}
}[
    - \meta{分块} : 用于在目录中显示的分块名称.
]

\newpage 作为示例, 下面给出部分\SimpleSystemTeX Manual的项目结构及其对应的导入命令:

\dirtree{%
    .0 .
    .1 {simplesystemtex.doc.tex} .
    .1 {sec/} .
    .2 {样式代码预设值.tex} .
    .2 {SimpleSystemTeX概述.tex} .
    .2 {SimpleSystemTeX命令索引.tex} .
    .2 {导言区/} .
    .3 {声明区块类型.tex} .
    .3 {样式设定.tex} .
    .2 {正文/} .
    .3 {创建区块.tex} .
    .3 {生成区块索引.tex} .
    .3 {添加链接.tex} .
    .2 {主文件/} .
    .3 {导入节文件.tex} .
    .3 {生成目录.tex} .
}

\Example{Manual项目结构-示例}{
    \par\noindent \cs{AddPart} \{\cs{SimpleSystemTeX}[bf]\}
    \par\noindent \cs{ImportSection} \{SimpleSystemTeX概述\}(\cs{SimpleSystemTeX}[bf]概述)
    \par\noindent \cs{AddPart} \{文件管理\}
    \par\noindent \cs{ImportSection} \{主文件/导入节文件\}
    \par\noindent \cs{ImportSection} \{主文件/生成目录\}
    \par\noindent \cs{AddPart} \{区块\cs{\&} 交叉引用\}
    \par\noindent \cs{ImportSection} \{导言区/声明区块类型\}
    \par\noindent \cs{ImportSection} \{正文/创建区块\}
    \par\noindent \cs{ImportSection} \{正文/生成区块索引\}
    \par\noindent \cs{ImportSection} \{正文/添加链接\}
    \par\noindent \cs{AddPart} \{其它\}
    \par\noindent \cs{ImportSection} \{导言区/样式设定\}
    \par\noindent \cs{AddPart} \{附录\}
    \par\noindent \cs{ImportSection} \{样式代码预设值\}
    \par\noindent \cs{ImportSection} \{SimpleSystemTeX命令索引\} (\cs{SimpleSystemTeX}[bf]命令索引)
}

\newpage
